\documentclass[12pt]{article}
\usepackage[utf8]{inputenc}
\usepackage[T1]{fontenc}
\usepackage{hyperref}
\title{Adasen Language Resources for K-12 Curriculum Development}
\author{Compiled by Claude Sonnet 5 (claude-sonnet-5) for Andr\'as Kornai}
\date{}
\begin{document}
\maketitle
\section{Introduction}
Adasen (ISO 639-3 tiu; Glottocode adas1235) is an Austronesian language of northeastern Abra and western Apayao, Cordillera, northern Luzon, spoken by one of the Itneg (Tinguian) subgroups. Sources found give a threatened status (EGIDS 6b), roughly 4,000 speakers out of about 10,300 ethnic Adasen (2010), and strong bilingualism with Ilocano; this is lower than the approximate figure of 10,000 in the reference description. Online documentation is thin. Note that Wikipedia's Itneg article associates Adasen with code itb, which conflicts with tiu used by other sources.
\section{Primary Educational Resources}
None found.
\section{Linguistic and Reference Materials}
The UP Diliman linguistics site has a language capsule on Adasen: \url{https://linguistics.upd.edu.ph/the-katig-collective/language-capsules/adasen/}. It reports that the only online work is a word list compiled by SIL in 1962--1966.
\section{Religious and Community Resources}
ScriptureEarth lists a New Testament (Ya Baru nga Tulag) with Global Recordings audio: \url{https://www.scriptureearth.org/00eng.php?iso=tiu}. People-group profile: \url{https://joshuaproject.net/people_groups/10155}.
\section{Digital and Audio Resources}
Audio is limited to the Global Recordings material linked via ScriptureEarth. Status page: \url{https://endangeredlanguages.com/node/87905}.
\section{Curriculum Development Context}
No Adasen school programme was found. Ilocano dominates; any curriculum would need new documentation work with the community, building on the word list and scripture text.
\end{document}